\documentclass[10pt]{article}
\usepackage[margin=0.85in]{geometry}
\usepackage{amsmath,amssymb,amsthm,hyperref}
\newtheorem{theorem}{Theorem}
\newtheorem{lemma}{Lemma}
\newcommand{\per}{\operatorname{per}}
\newcommand{\fall}[2]{#1^{\underline{#2}}}
\title{Two permanent congruences for shifted difference matrices}
\author{Principia Math in Progress}
\date{13 September 2026}
\begin{document}
\maketitle

\paragraph{Task 1. Prove both permanent congruences.}
Let \(D_N=[i-j]_{1\le i,j\le N}\) and \(B_N=I_N+D_N\).
Wang and Sun \cite{WS}, p.~3, conjecture
\[
 \per B_{p-1}\equiv3\pmod p,\qquad
 \per B_p\equiv1+4p\pmod{p^2}
\]
for odd primes \(p\). MathDB also includes \(p=2\) in the first
assertion; that case is immediate from \(B_1=(1)\).
We prove the stronger polynomial statements below. Congruences of
polynomials mean coefficientwise congruences.

\begin{theorem}\label{main}
For every odd prime \(p\),
\[
 \per(I_{p-1}+tD_{p-1})\equiv1+2t^{p-1}\pmod p,
 \qquad
 \per(I_p+tD_p)\equiv1+4pt^{p-1}\pmod{p^2}.
\]
\end{theorem}

We first prove a symmetric-polynomial fact. For variables
\(x_1,\ldots,x_N\), write
\[
 \per[\delta_{ij}+t(x_i-x_j)]
 =\sum_{d=0}^N c_d(x_1,\ldots,x_N)t^d,
 \qquad P_j=\sum_{i=1}^N x_i^j.
\]
The polynomial \(c_d\) is symmetric and homogeneous of degree \(d\).
Also \(c_0=1\), and \(c_d=0\) for odd \(d\), since transposition
replaces \(t\) by \(-t\).

\begin{lemma}\label{linear}
For \(1\le2m\le N\), the expression of \(c_{2m}\) in the power sums
has the form
\[
 c_{2m}=-(m-1)!\fall{N}{m}P_{2m}
    +\text{terms containing at least two positive-index power sums},
\]
where \(\fall{N}{m}=N(N-1)\cdots(N-m+1)\).
Every coefficient in this expression has denominator divisible only
by primes at most \(2m\).
\end{lemma}

\begin{proof}
The denominator assertion follows from the fundamental theorem of
symmetric polynomials and Newton's identities: an integral symmetric
polynomial of degree \(2m\) is integral in the elementary symmetric
polynomials, and converting those of degree at most \(2m\) to power
sums introduces only the integers \(1,\ldots,2m\) as divisors.
There is no constant term by homogeneity.

Here is an explicit calculation of the coefficient of \(P_{2m}\).
Define a linear functional \(\mathcal W\) on monomials in four
auxiliary variables by
\[
 \mathcal W(u^a r^b v^c s^d)=
 \begin{cases}a!c!,&a=b,\ c=d,\\0,&\text{otherwise}.\end{cases}
\]
Expansion of the rank-two factorization \(t(x_i-x_j)
=(tx_i,1)(1,-tx_j)^T\) gives
\begin{equation}\label{functional}
 \per[\delta_{ij}+t(x_i-x_j)]
 =\mathcal W\prod_{i=1}^N
 \bigl(1+vr+tx_i(ur-vs)-t^2x_i^2us\bigr).
\end{equation}
Indeed, selecting a subset of the factors selects a principal
permanent; the factorials count all matchings between the two
types of chosen row and column terms.

Put
\[
 \ell_j=[z^j]\log\bigl(1+vr+z(ur-vs)-z^2us\bigr),\quad j\ge1.
\]
The product in \eqref{functional} is
\[
 (1+vr)^N\exp\left(\sum_{j\ge1}t^jP_j\ell_j\right).
\]
Consequently the coefficient of \(P_{2m}\) in \(c_{2m}\) is
\(C_{2m}(N)=\mathcal W((1+vr)^N\ell_{2m})\).
These are formal coefficient calculations: only finitely many
auxiliary monomials can survive \(\mathcal W\) at any fixed power
of \(t\), so no analytic convergence is involved.

For clarity, expand the logarithm and let \(q\) be the number of
quadratic terms \(-z^2us\) chosen. The remaining \(2m-2q\) terms
come from \(zur\) and \(-zvs\). The functional forces us to take
exactly \(q\) factors \(vr\) from
\((1+vr)^{N-(2m-q)}\). Summing over the two linear choices gives
\[
 \sum_{a+b=2m-2q}(-1)^b
       \frac{(a+q)!}{a!}\frac{(b+q)!}{b!}
 =(q!)^2[z^{2m-2q}](1-z^2)^{-q-1}
 =(q!)^2\binom mq.
\]
Thus
\begin{equation}\label{finite}
 C_{2m}(N)=-\sum_{q=0}^m
 (2m-q-1)!q!\binom mq\binom{N-2m+q}{q}.
\end{equation}
The finite Chu--Vandermonde identity evaluates this sum as
\[
 C_{2m}(N)=-(2m-1)!
 \frac{(-N)_m}{(1-2m)_m}
 =-(m-1)!\fall{N}{m},
\]
where \((a)_m=a(a+1)\cdots(a+m-1)\).
For example, after dividing the sum in \eqref{finite} by
\(-(2m-1)!\), its summand is
\(\frac{(-m)_q(N-2m+1)_q}{(1-2m)_q q!}\), the terminating
binomial identity. None of its denominator factors is zero.
This proves the lemma.
\end{proof}

\begin{proof}[Proof of Theorem~\ref{main}]
First let \(N=p-1\) and reduce modulo \(p\).
Simultaneously permuting the row and column indices by multiplication
by any \(a\in\mathbb F_p^\times\) replaces \(t\) by \(at\).
Hence the polynomial \(\per(I+tD_N)\), of degree at most \(p-1\),
has no terms except degrees \(0\) and \(p-1\).
Its constant term is one. To calculate the top coefficient, for
arbitrary \(x_i\) the rank-two expansion gives
\[
 \per[x_i-x_j]=\sum_{k=0}^N
 (-1)^{N-k}k!(N-k)!e_k(x)e_{N-k}(x),
\]
where \(e_k\) denotes the elementary symmetric polynomial.
For \(x_i\) running through \(\mathbb F_p^\times\),
\(\prod_i(X-x_i)=X^{p-1}-1\). Thus only \(k=0,N\) remain,
\(e_N=-1\), and Wilson's theorem gives \(N!=-1\).
Since \(N\) is even, the two terms each equal one. This proves
the first polynomial congruence.

Now set \(N=p\), \(x_i=i\), and work modulo \(p^2\).
The usual finite-field power sums give
\[
 P_j\equiv0\pmod p\quad(1\le j<p-1),\qquad
 P_{p-1}\equiv-1\pmod p.
\]
For \(1\le2m\le p-1\), all denominators in Lemma~\ref{linear}
are invertible modulo \(p^2\). Every term with at least two power
sums vanishes modulo \(p^2\): each of its indices is less than
\(p-1\), and each factor is divisible by \(p\). Therefore
\begin{equation}\label{reduction}
 c_{2m}(1,\ldots,p)\equiv-(m-1)!\fall{p}{m}P_{2m}
 \pmod{p^2}.
\end{equation}
If \(2m<p-1\), both \(\fall p m\) and \(P_{2m}\) are divisible
by \(p\), so this coefficient is zero modulo \(p^2\).
For \(m=(p-1)/2\), we have
\[
 -(m-1)!\fall p m
 \equiv(-1)^m p\bigl((m-1)!\bigr)^2\pmod{p^2}.
\]
Wilson's theorem, paired around the midpoint, says
\((m!)^2\equiv(-1)^{m+1}\pmod p\); also
\(m^2\equiv1/4\pmod p\). It follows that
\[
 (-1)^m\bigl((m-1)!\bigr)^2\equiv-4\pmod p.
\]
Using \(P_{p-1}\equiv-1\pmod p\) in \eqref{reduction} gives
\(c_{p-1}\equiv4p\pmod{p^2}\).
The odd coefficients, including degree \(p\), vanish identically.
Together with \(c_0=1\), this proves the second polynomial
congruence. Specializing \(t=1\) proves both original assertions.
\end{proof}

\paragraph{Verification.}
The general proof above applies to every required prime. The accompanying
script checks the finite coefficient identity symbolically, isolates its
linear power-sum coefficient using finite-field roots of unity, and
compares complete permanent polynomials with the two polynomial
congruences for small primes. Independent direct permutation sums
also check the permanent algorithm at small orders.

\paragraph{Task 2. Lift the congruence.}
Determine the complete polynomial \(\per(I_p+tD_p)\) modulo \(p^3\),
including the contributions from products of two power sums.

\paragraph{Task 3. Vary the matrix order.}
Determine coefficientwise congruences modulo \(p^2\) for
\(\per(I_{kp}+tD_{kp})\) when \(k\ge2\), accounting for the
higher-degree terms whose power-sum expansions can have denominators
divisible by \(p\).

\paragraph{Task 4. Extend the kernel.}
Find analogous congruences for low-rank antisymmetric polynomial
kernels \(K(x,y)=-K(y,x)\).

\paragraph{Task 5. Evaluate exact permanents.}
Derive a closed formula or a short recurrence for
\(\per(I_N+tD_N)\) at arbitrary integer orders \(N\).

\begin{thebibliography}{1}
\bibitem{WS} H.~Wang and Z.-W.~Sun,
\emph{Evaluations of some Toeplitz-type determinants},
\href{https://arxiv.org/abs/2206.12317}{arXiv:2206.12317v8},
14 February 2023. The two conjectures appear together on p.~3
of the actual 22-page PDF.
\end{thebibliography}
\end{document}
